\section{Part-level routing}
\label{sec:routing}

\subsection{What the store already prunes, and what it does not}

\prior{} Predicate pruning in a MergeTree table happens at three levels: parts
excluded by the minimum and maximum of the partition key, granule ranges
selected by the sparse primary index, and granules excluded by skipping
indexes. All three read a predicate. Similarity has no predicate, so none of
them applies to it, and the vector index is consulted once per surviving part.
The same arrangement in Lucene and Milvus has the same
consequence~\cite{lucene,wang2021milvus}: the cost of a similarity query is a
sum over segments, and it grows as segments accumulate.

\subsection{Region summaries}

The bounds in this subsection are established \prior{}; their arrangement per
part is what is proposed \proposed{}. Rows inside a part are grouped into
regions --- by $k$-means over the dense vectors, or by granule ranges when the
ordering key already clusters similar rows. Region $R$ carries

\begin{itemize}[leftmargin=*,itemsep=1pt]
\item a centroid $\mu_R$ and a covering radius
  $r_R = \max_{x \in R}\lVert v(x) - \mu_R \rVert$;
\item a basis summary: $M_{R,i} = \max_{x \in R} c_i(x)$ for the elements
  present, truncated to the largest entries with one residual bound
  $M_{R,\ast}$ covering the rest;
\item statistics: the row count, the number of rows carrying each retained
  element, and quantiles of the score distribution of the region under a
  reference query model.
\end{itemize}

The centroid-and-radius pair is the routing object of metric trees --- the
branch-and-bound $k$-nearest-neighbour search of Fukunaga and
Narendra~\cite{fukunaga1975branch}, metric trees~\cite{uhlmann1991metric}, and
the M-tree's covering radius~\cite{ciaccia1997mtree} --- and bounding regions
for pruning go back to the R-tree~\cite{guttman1984rtree}. The statistics are
what distributed retrieval calls resource-selection evidence: CORI ranks
collections by term statistics~\cite{callan1995cori}, ReDDE estimates how many
relevant documents a collection holds~\cite{si2003redde}, Taily models each
shard's score distribution and selects on its tail~\cite{aly2013taily}, and
selective search over topically clustered shards is the arrangement in
which such selection pays~\cite{kulkarni2015selective}, resting on the older
cluster hypothesis~\cite{vanrijsbergen1979ir}. The proposal contributes the
container, not the bound or the statistic.

\begin{proposition}[Region bounds]
For non-negative sparse weights, every $x \in R$ satisfies
\[
  \sum_{i \in A(Q)} q_i\, c_i(x) \;\le\; \mathrm{UB}_s(Q,R)
  \;=\; \sum_{i \in A(Q)} q_i\, \widehat{M}_{R,i},
\]
where $\widehat{M}_{R,i}$ is $M_{R,i}$ when retained and $M_{R,\ast}$
otherwise. For the dense leg,
\[
  \langle u, v(x)\rangle \le \langle u, \mu_R\rangle + r_R \lVert u \rVert,
  \qquad
  \lVert u - v(x)\rVert \ge \lVert u - \mu_R \rVert - r_R .
\]
\end{proposition}
\begin{proof}
The sparse bound is termwise. The dense bounds are Cauchy--Schwarz and the
triangle inequality applied to $v(x) - \mu_R$, whose norm is at most $r_R$ by
the definition of the covering radius.
\end{proof}

\subsection{From regions to a table-level structure}

\proposed{} A part summary is the union of its region summaries: an enclosing
ball, the maximum of the region maxima per basis element, and summed
statistics. Part summaries are themselves grouped into a tree with a small
fan-out, each interior node carrying the same fields computed over its
children, which is the M-tree's construction applied to parts rather than to
rows. The structure is metadata: its size is proportional to the number of
regions, not to the number of rows, and it is maintained by the merge that
produces a part, since a merge replaces the summaries of its inputs with one
computed from the output it already has in memory.

\subsection{The planner rule}

\proposed{} A retrieval query is answered by best-first traversal over that
tree. The planner keeps a priority queue of nodes ordered by
$\mathrm{UB}(Q,\cdot)$ and a running threshold $\theta$, the $k$-th best score
found. It pops the node of highest bound; if the bound is at most $\theta$ it
stops, because no descendant can beat the threshold. Interior nodes expand to
children, parts expand to their regions, and a region is evaluated by the
posting lists and, for the dense leg, by the part's approximate graph
restricted to the region. This is branch and bound over the same bounds the
metric-tree literature uses, with the threshold rule of Fagin et
al.~\cite{fagin2003ta}.

Two departures from exactness are deliberate and are reported separately. The
per-part graph search is approximate, so the answer is exact only relative to
what that search returns; recall is governed by its parameters, and the
protocol fixes recall rather than the parameters. And an approximate routing
variant stops when the best remaining bound falls below $\theta/(1+\varepsilon)$,
or after a budget of parts, which is the same trade SPANN makes when it prunes
posting lists by their centroid distance relative to the
closest~\cite{chen2021spann}.

Predicates compose with routing rather than replacing it: partition pruning and
primary-index ranges narrow the parts before summaries are read, and a
selective predicate can make routing unnecessary. The hard case for every
system in this family is a predicate that is selective and uncorrelated with
the index structure, which is what filtered graph search
addresses~\cite{gollapudi2023filtered,patel2024acorn}; the design as written
carries the predicate into the region scan and does not solve the general case.

\subsection{Where the bounds stop working}

\proposed{} \emph{Interpretation, not measurement.} Three failure modes are
predictable and the protocol measures each.

\textbf{High intrinsic dimension.} Covering radii of $k$-means regions approach
the distances between centroids, every region's ball bound exceeds $\theta$,
and the traversal visits everything. That is the classical result that
partition-based structures degrade to a scan in high
dimension~\cite{weber1998quantitative,beyer1999nn}; the routed store then costs
the baseline plus the summaries it read, and the protocol reports the summary
cost separately so the overhead is visible rather than folded into a total.

\textbf{Heavy basis elements.} An element carried with high weight by rows
everywhere raises $\mathrm{UB}_s$ for every region at once. Dynamic pruning has
the same weakness for common terms, which is why impact-ordered and
score-at-a-time evaluation exist~\cite{anh2002impact,lin2015anytime}.

\textbf{Small parts.} A freshly inserted part holds few rows and a summary
computed over them, so its bound is tight but its share of the corpus is
small; many such parts make the summary scan the dominant term until merges
consolidate them. Part count is therefore a controlled variable of the
protocol, not a detail of the deployment.

\subsection{Routing by summaries elsewhere}

\prior{} Summary-first retrieval appears in the retrieval-augmented generation
literature in a different form. GraphRAG partitions an entity graph into
communities, summarises each with a language model, and answers a global query
by mapping over community summaries and reducing their
partials~\cite{edge2024graphrag}; RAPTOR builds a tree of recursive abstractive
summaries and retrieves from its levels~\cite{sarthi2024raptor}. The structural
analogy is exact --- a summary per region, consulted before the rows --- and
the mechanism is not: those summaries are natural language read by a model, and
the map step visits every community at a level rather than excluding any by a
bound. The routing here excludes regions arithmetically and reads no summary
text.
